% Pursuing Stacks (cotes 134-2 to 134-8): archive cote 134-5, site page 64 (PDF page 65), typescript page 251.
% typescript. Draft compared: ps4.tex:3215-3280.
% AI-transcribed by Claude Opus 5.5 (claude-opus-5-5, adaptive thinking, effort high), Message Batches API, batch ps2 (msgbatch_01NVwr1Yuo8qJf5rykeGBXPK), 2026-11-12; not yet checked by eye.
\scanpage{134-5/64}
\mspage{251}

Let
\[ h \;=\; h_{\underline{W}_A} \]
this homotopy structure, hence a corresponding notion of h-\uline{equivalence} \inserted[typed]{\hand{(}or h-\uline{homotopy}}\hand{)}
$\overset{h}{\hand{\sim}}$ for arrows in $A\hat{}$, a corresponding notion of h-homotopisms i.e.
a set of arrows
\[ W^h \subset \mathrm{Fl}\,A\hat{}\ \hand{,} \quad \text{such that} \quad W^h \subset \underline{W}_A \quad , \]
a notion  of h-homotopy interval (namely an interval such that $\hand{\delta}_0$ and $\hand{\delta}_1$
be  h-homotopic , for which it is sufficient\hand{,}but not necessary\hand{,}that I
be locally aspheric ...), and last not least, a notion of \uline{contractible}
\uline{objects}, making up a full subcategory
\[ A\hat{}_c \subset A\hat{}\ , \quad \text{such that} \quad A\hat{}_c \subset A\hat{}_{\mathrm{loc.as}} \quad . \]
The latter inclusion, in case A is aspheric, can be equally written
\[ A\hat{}_c \subset A\hat{}_{\mathrm{tot.as}} \quad \text{(if A aspheric).} \]

Coming back now to the contractibility structure $(M,M_c)$\hand{,} and a
functor  i: A $\longrightarrow$ M , we are interested in the corresponding functor
\[ u = i^* : M \longrightarrow \xout{\illegible{1}}\ M' \;=\; A\hat{}\ \hand{,} \]
where both members will be viewed as being endowed with their re\unsure{sp}ective
homotopy structures  -  the one of M being of the most restrictive type
envisioned in section 51, namely it is  defined in terms of a contracti-%
bility structure, whereas the one of \deleted{xx}\inserted[typed]{M'}is a priori defined in terms of
a homotopy interval structure, but not necessarily in terms of a contrac
tibility structure. Now this situation has been described in section
5\unsure{3} , as far as compatibility \inserted[typed]{conditions}with homotopy structures are concerned ,
independently by the way of any special assumption on M'  (such as being
a category of presheaves on some small category A), or on the functor
u, except for commuting to finite products. Compatibility of u \xout{xxx} with
the homotopy structures on M,M'  can be expressed by either one of the
following six conditions, which are all equivalent:

H 1)  u transforms homotopic arrows of M into homotopic arrows in M'\hand{.}

H 2)  u transforms homotopisms \deleted{\illegible{1}}\inserted{in} M into homotopisms in M' .

H 3)  u transforms any homotopy interval $\underline{I} = (I,\hand{\delta}_0,\hand{\delta}_1)$ in M into a
homotopy interval in M' \hand{(} i.e. if two sections of an object
I over $e_M$ are homotopic, so are $u(\hand{\delta}_0)$ and $u(\hand{\delta}_1)$\hand{).}

H 3')  Same as H 3), but $\underline{I}$ being restricted\hand{,} to be in a given \deleted{genera-}%
\deleted{ting} family of homotopy intervals\hand{,}\inserted{generating}for the homotopy \xout{intervals}
structure in M.

\margin{64}
